\chapter{Bedwetting in Older Children}

\section*{Definition and Overview}
Bedwetting in older children, also called nocturnal enuresis, is the repeated involuntary passing of urine during sleep in a child aged about five years or older. It is common and usually not caused by a physical or emotional problem. Most children become dry at night as their bladder control matures, but some continue to wet the bed for years. Bedwetting is not the child's fault, and with support, reassurance, and simple treatments most children can become dry. Daytime wetting or other symptoms point to possible underlying causes.

\section*{Epidemiology}
Nocturnal enuresis is very common in childhood. Roughly 1 in 7 five-year-olds, 1 in 12 seven-year-olds, and about 1 in 50 fifteen-year-olds still wet the bed. It is slightly more common in boys than girls. A family history is very common: children whose parents were bed-wetters are much more likely to be affected. The condition gradually improves with age, and most children outgrow it without treatment. It is not linked to intelligence, behaviour problems, or poor parenting.

\section*{Aetiology and Pathophysiology}
\begin{itemize}
    \item \textbf{Delayed bladder control:} the bladder and its nerve pathways mature at different rates in different children
    \item \textbf{High overnight urine production:} some children produce too much urine at night because the hormone (vasopressin) that normally concentrates urine overnight is not yet produced in sufficient amounts
    \item \textbf{Small functional bladder capacity:} the bladder holds less urine than expected for the child's age
    \item \textbf{Deep sleep:} children who sleep very deeply may not wake when the bladder is full
    \item \textbf{Constipation:} a full bowel can press on the bladder and reduce its capacity
    \item \textbf{Family history:} a strong genetic component; about three-quarters of affected children have a parent who was also a bed-wetter
\end{itemize}

\section*{Clinical Features}
\begin{itemize}
    \item Wetting the bed during sleep at least once a month, in a child aged five or older
    \item The child may be dry for some nights and wet on others, in any pattern
    \item Usually the child has no daytime symptoms and is otherwise healthy
    \item Some children also have daytime urgency, frequency, or occasional daytime wetting
    \item Associated symptoms such as persistent daytime wetting, painful urination, or thirst may indicate another cause
    \item Distress, embarrassment, or avoidance of sleepovers and school camps
\end{itemize}

\section*{Differential Diagnosis}
\begin{itemize}
    \item \textbf{Urinary tract infection:} painful or frequent urination, and daytime wetting
    \item \textbf{Diabetes (type 1 or 2):} excessive thirst, weight loss, and increased urination
    \item \textbf{Constipation:} a common and often overlooked cause
    \item \textbf{Bladder or kidney disorders,} including overactive bladder or structural abnormalities
    \item \textbf{Obstructive sleep apnoea:} loud snoring and daytime sleepiness may be linked to bedwetting
    \item \textbf{Stress or emotional factors,} though these are rarely the sole cause
    \item \textbf{Side effects of medicines,} such as some asthma treatments or diuretics
\end{itemize}

\section*{When to Seek Medical Care}
See a doctor if bedwetting starts suddenly after the child has been dry for six months or more, if there is daytime wetting or painful urination, or if the child is seven years or older and you want treatment options. A single consult is also worthwhile if the wetting is frequent or distressing.

\textbf{Call 000 (or local emergency services) for:}
\begin{itemize}
    \item Painful urination, fever, or blood in the urine
    \item Excessive thirst, weight loss, or extreme tiredness
    \item Pain in the abdomen, back, or genitals
    \item Sudden severe symptoms or the child appearing unwell
\end{itemize}

\section*{Diagnosis and Investigations}
\begin{itemize}
    \item \textbf{History:} whether wetting is only at night or also during the day, the pattern of wet nights, fluid intake, bowel habit (constipation), and family history
    \item \textbf{Bladder diary:} recording wet and dry nights, daytime toilet visits, and drinks over a week or two
    \item \textbf{Urine test:} a dipstick test to exclude urinary tract infection, diabetes, or kidney problems
    \item \textbf{Examination:} checking the abdomen for a full bladder or constipation, and examining the genitals
    \item \textbf{Ultrasound:} occasionally used to assess bladder size and emptying, or to check the kidneys
    \item \textbf{Further tests:} only ordered if specific concerns arise, such as blood tests for diabetes
\end{itemize}

\section*{Management}
\begin{itemize}
    \item \textbf{Reassurance and education:} explain to the child and parents that bedwetting is not anyone's fault and usually resolves with time
    \item \textbf{Regular toileting:} encourage emptying the bladder at bedtime and adequate daytime toileting
    \item \textbf{Bladder training:} for some children, exercises and timing routines to improve bladder capacity and control
    \item \textbf{Constipation treatment:} treating constipation if present, which often improves wetting
    \item \textbf{Enuresis alarm:} a device that wakes the child when wetting begins, training the child to wake at bladder fullness; it is the most effective treatment and requires a motivated child and family
    \item \textbf{Medication:} desmopressin (a tablet or melt that reduces overnight urine production) may be used, especially for short-term situations such as sleepovers; other medicines may be used for overactive bladder
    \item \textbf{Support:} education and regular follow-up; reward effort, not dryness, to keep motivation positive
\end{itemize}

\section*{Complications}
\begin{itemize}
    \item Low self-esteem, embarrassment, and social withdrawal
    \item Avoidance of sleepovers, camps, and school activities
    \item Family stress, conflict, and blame placed on the child
    \item Sleep disturbance for the child and household
    \item Rarely, skin irritation from prolonged contact with wet bedding
    \item Underlying conditions such as infection or diabetes may go unnoticed if the wetting is dismissed
\end{itemize}

\section*{Prognosis and Outcomes}
Most children eventually become dry at night without treatment; the condition resolves spontaneously at a rate of about 1 in 7 affected children each year. With an enuresis alarm, around two-thirds to three-quarters of children become dry or significantly improve within a few months, though some relapse. Medication can provide rapid but often temporary dryness. Bedwetting that persists into adolescence is more persistent, but even then most young people eventually achieve dryness.

\section*{Prevention and Self-Care}
\begin{itemize}
    \item Do not punish or blame the child; treat bedwetting as a developmental stage, not a behaviour problem
    \item Limit large drinks in the hour or two before bed, while ensuring the child drinks enough through the day
    \item Encourage a regular toilet visit at bedtime
    \item Avoid drinks containing caffeine before bed
    \item Treat constipation early with diet, fluids, and activity
    \item Use a waterproof mattress protector and involve the child in changing the bedding without shame
    \item Reward the child for following routines and for dry nights, praising effort rather than pressuring for results
    \item Seek advice early if daytime wetting, pain, or new symptoms develop
\end{itemize}

\section*{Sources and Further Reading}
The following sources provide reliable, up-to-date information on bedwetting in children.
\begin{itemize}
    \item NHS. \textit{Bedwetting in children.} https://www.nhs.uk/conditions/
    \item NICE. \textit{Bedwetting in children and young people.} https://www.nice.org.uk
    \item Mayo Clinic. \textit{Bed-wetting.} https://www.mayoclinic.org
    \item MedlinePlus. \textit{Bedwetting.} https://medlineplus.gov
    \item MSD Manual. \textit{Nocturnal Enuresis (Bedwetting).} https://www.msdmanuals.com
\end{itemize}
